% Fichier généré automatiquement par tools/generate_corrections.py.
% Source : PERF/PERF-02-Marges/61_Hemostase/61_Hemostase.tex
% Statut : partiel (5/7 question(s) renseignée(s), 0 reprise(s)).
\begin{corrigebox}[Corrigé partiel extrait de la banque]
\CorrigeQuestion{1}
$H_{\text{bo}}(p) = H_{\text{cor}}(p) \dfrac{K_1}{p\left(1+T_m p\right)}$ $=  \dfrac{K_1 K_p}{p\left(1+T_m p\right)}$.

$H_{\text{cr}}(p)= -K_2\dfrac{\dfrac{K_1}{p\left(1+T_m p\right)}}{1+H_{\text{cor}}(p) \dfrac{K_1}{p\left(1+T_m p\right)}}$

$= -K_2\dfrac{K_1}{p\left(1+T_m p\right)+H_{\text{cor}}(p) K_1}$
$= -\dfrac{K_1K_2}{p\left(1+T_m p\right)+K_p K_1}$
\CorrigeQuestion{2}
Exprimons $\varepsilon(p)$ en fonction de $Z_c(p)$ et $C_{r}(p)$ :

$
\varepsilon(p)=Z_c(p)-Z(p) = Z_c(p)- \left(\varepsilon(p) H_{\text{cor}}(p)-   K_2 C_r(p)\right) \dfrac{K_1}{p\left(1+T_m p\right)}$

$\Leftrightarrow \varepsilon(p)\left(1 +  H_{\text{cor}}(p)\dfrac{K_1}{p\left(1+T_m p\right)} \right) 
= Z_c(p) + K_2 C_r(p) \dfrac{K_1}{p\left(1+T_m p\right)}$

$\Leftrightarrow \varepsilon(p)  = 
Z_c(p)\dfrac{1}{1 +  H_{\text{cor}}(p)\dfrac{K_1}{p\left(1+T_m p\right)}} 
+   K_2 C_r(p) \dfrac{K_1}{p\left(1+T_m p\right)} \dfrac{1}{1 +  H_{\text{cor}}(p)\dfrac{K_1}{p\left(1+T_m p\right)}}$

$\Leftrightarrow \varepsilon(p)  = 
Z_c(p)\dfrac{{p\left(1+T_m p\right)}}{{p\left(1+T_m p\right)} +  H_{\text{cor}}(p){K_1}} 
+  K_2 C_r(p)  \dfrac{K_1}{p\left(1+T_m p\right) +  H_{\text{cor}}(p){K_1}}$

En prenant une entrée échelon et une perturbation échelons, on a $Z_c(p) = \dfrac{Z_{c0}}{p}$ et 
$C_{r}(p) = \dfrac{C_{r0}}{p}$.

On a donc $\lim\limits_{t\to +\infty} \varepsilon(t) $
$=\lim\limits_{p\to 0} p\varepsilon(p)$
$=  \lim\limits_{p\to 0} Z_{c0}\dfrac{{p\left(1+T_m p\right)}}{{p\left(1+T_m p\right)} +  H_{\text{cor}}(p){K_1}} 
+   K_2 C_{r0}  \dfrac{K_1}{p\left(1+T_m p\right) +  H_{\text{cor}}(p){K_1}} $ 
$ =   \dfrac{K_2 C_{r0}}{ K_p}$.

AN : $\varepsilon_s < \SI{1}{mm}$
$\Leftrightarrow \dfrac{K_2 C_{r0}}{ K_p} < \SI{1}{mm}$  
$\Leftrightarrow 
2,78 \cdot 10^{-2} \times 2,7 \cdot 10^{-3} \times 10^3<  K_p$  soit $K_p >0,08$.
\CorrigeQuestion{3}
En ajoutant le gain de 0,08, il faut translater la courbe de gain vers le bas de 22 dB.

\begin{marginfigure}
\includegraphics[width=\linewidth]{cor_06}
\end{marginfigure}

La marge de phase est supérieure à 60\degres.
\CorrigeQuestion{4}
\CorrigeACompleter{Aucun élément de réponse n'est présent dans la source d'origine pour cette question.}
\CorrigeQuestion{5}
\CorrigeACompleter{Aucun élément de réponse n'est présent dans la source d'origine pour cette question.}
\CorrigeQuestion{6}
\begin{center}
\begin{tabular}{|l|c|c|c|}
\hline
 & P & PI & PID \\  \hline
 Temps de réponse < à 5\,\% < \SI{0,2}{s} &Ok & Ok & Ok \\ \hline
Précision < \SI{1}{mm} & Ok (?) & Ok & Ok \\ \hline
 Dépassement < à 10\,\% < \SI{0,2}{s} & Pas Ok & Pas Ok & Ok \\ \hline
\end{tabular}
\end{center}

\begin{marginfigure}
\includegraphics[width=\linewidth]{cor_07}
\end{marginfigure}
\CorrigeQuestion{7}
\begin{marginfigure}
\includegraphics[width=\linewidth]{cor_08}
\end{marginfigure}

La marge de phase est supérieure à 60\degres.
\end{corrigebox}
